\documentclass[11pt]{article}

% ---------------------------------------------------------------------------
% Design note: the mathematics of a buprenorphine mu-receptor interaction
% model for stanpumpR.  Source of truth for docs/buprenorphine-interaction.pdf.
% Compile:  pdflatex buprenorphine-interaction.tex   (run twice for refs)
% ---------------------------------------------------------------------------

\usepackage[a4paper,margin=1in]{geometry}
\usepackage{amsmath,amssymb}
\usepackage{booktabs}
\usepackage{enumitem}
\usepackage{xcolor}
\usepackage[hidelinks]{hyperref}

\setlength{\parskip}{0.4em}
\setlength{\parindent}{0pt}
\setlength{\emergencystretch}{2.5em}

% A lightweight "correction" callout: a coloured run-in label plus a thin rule,
% using no framing package, so the document compiles against a minimal TeX Live.
\definecolor{corrcol}{RGB}{140,20,20}
\newenvironment{correction}[1]
  {\par\vspace{0.3em}\noindent\textcolor{corrcol}{\rule{\linewidth}{0.4pt}}\par
   \noindent\textbf{\textcolor{corrcol}{[Correction] #1}}\par\nopagebreak\vspace{0.2em}}
  {\par\noindent\textcolor{corrcol}{\rule{\linewidth}{0.4pt}}\par\vspace{0.3em}}

\newcommand{\dd}[2]{\frac{\mathrm{d}#1}{\mathrm{d}#2}}
\newcommand{\Ce}{C_{\mathrm e}}
\newcommand{\Cesub}[1]{C_{\mathrm{e},#1}}
\newcommand{\kon}[1][]{k_{\mathrm{on}#1}}
\newcommand{\koff}[1][]{k_{\mathrm{off}#1}}
\newcommand{\rhoz}{\rho_0}

\title{\textbf{The mathematics of a buprenorphine $\mu$-receptor interaction model}\\[0.3em]
\large A design note for stanpumpR}
\author{Prepared with Claude Code at the request of Steven~L.\ Shafer}
\date{2026-11-19 \quad\textbf{Status: proposal, not implemented}}

\begin{document}
\maketitle

\begin{abstract}
\noindent
This note records the mathematics a future buprenorphine interaction engine
would rest on, so the derivations are settled before any code is written.
Nothing here changes the app. It began as an assessment of an external problem
brief (written for ChatGPT); the structure and most of the results in that brief
are correct and are adopted. Four places where this note \emph{departs from or
strengthens} the brief are flagged with \textcolor{corrcol}{\textbf{[Correction]}}
callouts: P2 is linear (indeed compartmental), not nonlinear; P4 has an exact
invariant, not merely a bounded deviation; P1's closed form is exact but not
automatically best-conditioned; and P3's reduction to Olofsen's linear
transduction requires a transduction Hill coefficient of one.
\end{abstract}

% ===========================================================================
\section{Scope and the kinetic premise}

\textbf{$\mu$ receptor only.} Ligands: buprenorphine; the full agonists
fentanyl and sufentanil; the antagonists naloxone, naltrexone, nalmefene.
$\kappa/\delta$, ORL-1/NOP and norbuprenorphine are out of scope.

The model is a \textbf{pharmacodynamic post-process} on the effect-site
concentration curves $\Cesub{i}(t)$ the simulator already produces. It must not
touch the pharmacokinetics, and it must reduce to the existing additive opioid
scale $U=\sum_j \Cesub{j}/\mathrm{MEAC}_j$ when no partial agonist or antagonist is
present (\S\ref{sec:p4}).

The premise that forces a \emph{kinetic} rather than an equilibrium model:
buprenorphine's clinical signature --- durable blockade of other opioids,
precipitated withdrawal, the ventilatory ceiling --- is set by its \textbf{slow
dissociation} ($\koff[,B]\approx 0.019~\mathrm{min^{-1}}$,
$t_{1/2}\approx 37$~min), not by its affinity. A full agonist such as
sufentanil binds with comparable or higher affinity yet is readily displaced,
because its dissociation is fast. Displaceability is governed by the
dissociation rate constant / residence time, which \textbf{cancels out of an
equilibrium description}: at equilibrium only $K_D=\koff/\kon$ survives, and
$K_D$ does not distinguish a fast ligand from a slow one of the same affinity.
Carrying $\koff$ explicitly is the whole point of the engine.

% ===========================================================================
\section{The model (fixed published structure)}
\label{sec:model}

Index the $\mu$ ligands by $i$. Each has an effect-site concentration
$\Cesub{i}(t)$ and rate constants $\kon[,i],\koff[,i]$
($K_{D,i}=\koff[,i]/\kon[,i]$). Normalise the receptor pool to~1; let $\rho_i$
be the fraction bound by ligand $i$ and $\rhoz=1-\sum_j\rho_j$ the free
fraction. Binding kinetics (Olofsen et al.\ 2022; Yassen et al.\ 2006/2007):
%
\begin{equation}
\dd{\rho_i}{t}=\kon[,i]\,\Cesub{i}(t)\,\rhoz(t)-\koff[,i]\,\rho_i(t),
\qquad \rhoz=1-\sum_j\rho_j .
\label{eq:bind}
\end{equation}
%
Units: $\kon$ in $\mathrm{mL\,ng^{-1}\,min^{-1}}$, $\Ce$ in $\mathrm{ng\,mL^{-1}}$,
so $\kon\Ce$ is $\mathrm{min^{-1}}$, as is $\koff$; $C_{50}=\koff/\kon$ is in
$\mathrm{ng\,mL^{-1}}$.

Ligands split by residence time $\tau=1/\koff$:
\begin{itemize}[nosep,leftmargin=1.4em]
\item \textbf{Slow} --- buprenorphine ($\koff\approx0.019$) and the
long-residence antagonists naltrexone, nalmefene --- keep a differential
state~\eqref{eq:bind}.
\item \textbf{Fast} --- fentanyl, sufentanil, naloxone --- reach their
\emph{conditional equilibrium} within whatever pool the slow ligands leave free.
Setting $\dd{\rho_j}{t}=0$ gives $\rho_j=\phi_j\rhoz$, $\phi_j=\Cesub{j}/C_{50,j}$;
summing $\rhoz+\sum_{\mathrm{slow}}\rho+\rhoz\sum\phi=1$:
\end{itemize}
%
\begin{equation}
\rho_j=\bigl(1-\textstyle\sum_{\mathrm{slow}}\rho\bigr)\,
\frac{\phi_j}{1+\sum_k\phi_k},\qquad
\phi_j=\frac{\Cesub{j}}{C_{50,j}},\quad C_{50,j}=\frac{\koff[,j]}{\kon[,j]}.
\label{eq:qss}
\end{equation}
%
Define the \textbf{crowding factor}
$\Phi(t)=\sum_{\mathrm{fast}}\phi_k(t)=\sum \Cesub{k}/C_{50,k}$. Then the total
fast-bound fraction is $\bigl(1-\sum_{\mathrm{slow}}\rho\bigr)\Phi/(1+\Phi)$ and
%
\begin{equation}
\rhoz=\frac{1-\sum_{\mathrm{slow}}\rho}{1+\Phi}.
\label{eq:rho0}
\end{equation}
%
\textbf{Effect (ventilation), linear transduction (Olofsen):}
%
\begin{equation}
E=E_0\Bigl(1-\sum_i\alpha_i\rho_i\Bigr),
\label{eq:effect}
\end{equation}
%
$\alpha=1$ for a full agonist (no ceiling), $\alpha<1$ for buprenorphine (the
ceiling), $\alpha=0$ for a pure antagonist. Between grid points each $\Cesub{i}(t)$
is \textbf{piecewise-linear in $t$}; the engine advances occupancy exactly
across one such interval given the endpoint concentrations.

\begin{table}[h]
\centering
\caption{Fitted parameters (Olofsen et al.\ 2022, ventilation).}
\begin{tabular}{lcc}
\toprule
Parameter & Opioid-naive & Chronic high-dose \\
\midrule
$\kon[,B]$ (mL\,ng$^{-1}$\,min$^{-1}$) & 0.26 & 0.085 \\
$\koff[,B]$ (min$^{-1}$; $t_{1/2}\approx37$~min) & 0.019 & 0.019 \\
$\alpha_B$ & 0.81 & 0.48 \\
$C_{50,F}=\koff[,F]/\kon[,F]$ (ng\,mL$^{-1}$) & 0.68 & 2.23 \\
$\alpha_F$ (fixed) & 1 & 1 \\
$E_0=V_B$ (L\,min$^{-1}$) & 20.2 & 20.2 \\
\bottomrule
\end{tabular}
\end{table}

% ===========================================================================
\section{One grid interval: notation}
\label{sec:interval}

On $[0,h]$ write every input as linear in $t$:
%
\begin{align}
\kon[,B]\,B(t)&=r+s\,t, & r&=\kon[,B]B_0, & s&=\kon[,B](B_1-B_0)/h,\\
1+\Phi(t)&=p+q\,t, & p&=1+\Phi_0, & q&=(\Phi_1-\Phi_0)/h,
\end{align}
%
with $B_0,B_1$ the buprenorphine endpoint concentrations and $\Phi_0,\Phi_1$ the
crowding factor at the two ends. $p\ge 1>0$ throughout (a crowding factor is
non-negative), which keeps every denominator and logarithm below well defined.

% ===========================================================================
\section{P1 --- the single-slow-ligand scalar ODE}

With buprenorphine the only slow ligand, $\sum_{\mathrm{slow}}\rho=x\equiv
\rho_{\mathrm{bup}}$, so $\rhoz=(1-x)/(1+\Phi)$. Substituting into
\eqref{eq:bind}:
%
\begin{equation}
\dd{x}{t}=a(t)\,(1-x)-\koff[,B]\,x,\qquad
a(t)=\frac{\kon[,B]B(t)}{1+\Phi(t)} .
\label{eq:scalar}
\end{equation}
%
This matches the brief. It is \textbf{linear in $x$} with a time-varying
coefficient,
%
\begin{equation}
\dd{x}{t}+\bigl[a(t)+\koff[,B]\bigr]x=a(t),
\label{eq:scalarlin}
\end{equation}
%
and with the \S\ref{sec:interval} parametrisation $a(t)=(r+st)/(p+qt)$ is a
ratio of two linear functions (M\"obius). Polynomial division gives the form
used throughout:
%
\begin{equation}
a(t)=A+\frac{w}{p+q\,t},\qquad A=\frac{s}{q},\qquad
w=r-\frac{sp}{q}=\frac{rq-sp}{q}.
\label{eq:adecomp}
\end{equation}
%
(When $q=0$, i.e.\ $\Phi$ constant on the step, use the limit in
\S\ref{sec:p1deg}.)

% ===========================================================================
\section{P1 --- exact advance, and its numerically safe evaluation}

\subsection{Closed form via the integrating factor}
\label{sec:p1closed}

Let $\kappa=\koff[,B]$. The integrating-factor exponent is
%
\begin{equation}
\Psi(t)=\int_0^t\!\bigl[a(u)+\kappa\bigr]\,\mathrm{d}u
=(A+\kappa)\,t+m\,\ln\!\frac{p+q\,t}{p},\qquad m=\frac{w}{q},
\label{eq:psi}
\end{equation}
%
so $\mu(t)=e^{\Psi(t)}=e^{(A+\kappa)t}\bigl((p+qt)/p\bigr)^{m}$ with
\textbf{non-integer exponent $m$} --- the source of the special function below.
The solution of \eqref{eq:scalarlin} is
%
\begin{equation}
x(h)=e^{-\Psi(h)}\Bigl[\,x(0)+\int_0^h e^{\Psi(t)}a(t)\,\mathrm{d}t\,\Bigr].
\label{eq:p1sol}
\end{equation}
%
Using \eqref{eq:adecomp} and substituting $y=p+qt$, the source integral reduces
to two integrals of the form $\int e^{\beta y}y^{c}\,\mathrm{d}y$, each a
\textbf{lower incomplete gamma function} $\gamma(c+1,\cdot)$ (an exponential
integral when $c$ is a negative integer):
%
\begin{equation}
\int_0^h e^{\Psi}a\,\mathrm{d}t=A\,J(m)+w\,J(m-1),\qquad
J(c)=\int_0^h e^{(A+\kappa)t}\Bigl(\tfrac{p+qt}{p}\Bigr)^{c}\mathrm{d}t .
\end{equation}
%
So the advance is \textbf{closed-form, with a logarithm (in $\Psi$) and an
incomplete-gamma term (in the source)} --- exactly the brief's
characterisation.

\subsection{Overflow-proof evaluation (do this, not \S\ref{sec:p1closed} directly)}
\label{sec:p1overflow}

$\mu(t)$ can overflow, and $\bigl((p+qt)/p\bigr)^{m}$ alone can overflow or
underflow while the product in \eqref{eq:p1sol} is finite. \textbf{Never form
$\mu$ or $e^{+\Psi}$ explicitly.} Write the advance as a weighted average with a
\emph{bounded} kernel:
%
\begin{equation}
x(h)=G(0,h)\,x(0)+\int_0^h G(t,h)\,a(t)\,\mathrm{d}t,
\label{eq:p1stable}
\end{equation}
\begin{equation}
G(t,h)=\exp\!\Bigl(-\!\int_t^h[a(u)+\kappa]\,\mathrm{d}u\Bigr)
=\exp\!\Bigl(-(A+\kappa)(h-t)-m\ln\tfrac{p+qh}{p+qt}\Bigr).
\label{eq:kernel}
\end{equation}
%
Because $a\ge 0$ and $\kappa>0$, the inner integral $\int_t^h(a+\kappa)\ge 0$,
so $G(t,h)\in(0,1]$ for every $t$. \textbf{Evaluate the combined exponent as one
number} --- do not compute $e^{(\cdot)}\cdot\mathrm{pow}(\cdot)$ as separate
factors, because the $\mathrm{pow}$ factor alone can overflow even though
$G\le1$.

\subsection{Degenerate limits}
\label{sec:p1deg}

\begin{itemize}[nosep,leftmargin=1.4em]
\item \textbf{$a\equiv0$} ($B=0$): pure decay, $x(h)=x(0)\,e^{-\kappa h}$.
\item \textbf{Constant $a$} ($s=0$ and $q=0$, or $w=0$): constant-coefficient
case, $x(h)=x_{\mathrm{ss}}+\bigl(x(0)-x_{\mathrm{ss}}\bigr)e^{-(a+\kappa)h}$,
$x_{\mathrm{ss}}=a/(a+\kappa)$. The $w=0$ branch ($rq=sp$: numerator and
denominator of $a$ share a root --- the brief's ``coincident denominators'')
collapses to this; test $|w|$ against a tolerance.
\item \textbf{$q\to0$} ($\Phi$ constant, $a$ linear in $t$): $m=w/q$ diverges
but $m\ln\bigl((p+qt)/p\bigr)\to w\,t/p$ is finite. Implement the limit via the
series $\ln(1+qt/p)/q\to t/p$, \emph{not} the formula \eqref{eq:psi}, or the
$m\ln$ term cancels catastrophically. Here $\Psi$ is quadratic in $t$ and the
source integral is an \textbf{error-function} term rather than incomplete-gamma.
\item \textbf{General guard}: branch on $|q|$ and $|w|$ small before evaluating
\eqref{eq:adecomp}--\eqref{eq:psi}.
\end{itemize}

\subsection{Conditioning and sub-stepping}

\begin{itemize}[nosep,leftmargin=1.4em]
\item \textbf{Large $(A+\kappa)h$} (many time constants per step): $G(0,h)$
underflows and $x\to$ quasi-steady; correct, but the source quadrature needs
nodes concentrated near $t=h$ where $G=O(1)$. Sub-step so that $\int_t^h(a+\kappa)$
per sub-step is $O(1)$ and $\Phi$ varies little across it.
\item \textbf{Large $|m|$ with small $q$}: use the \S\ref{sec:p1deg} series.
\end{itemize}

\begin{correction}{Exact is not the same as best-conditioned}
The incomplete-gamma form of \S\ref{sec:p1closed} is exact, but its evaluation
at \emph{non-integer, possibly negative} order $m$ has its own cancellation
regions. Because $a(t)$ is smooth and bounded on one grid step, a 3--4 node
\textbf{Gauss--Legendre quadrature of the bounded kernel}
in~\eqref{eq:p1stable} reaches machine precision with none of that fragility.
The recommendation: keep the closed form of \S\ref{sec:p1closed} as the
\emph{reference / validator}, and use the bounded-kernel quadrature
of~\eqref{eq:p1stable} as the production path. ``Load-bearing'' describes the
\emph{exactness requirement}, which \eqref{eq:p1stable} also meets; it need not
mean the special-function evaluation ships.
\end{correction}

\subsection{Validation}
Check $x(h)$ against an adaptive high-order integration of \eqref{eq:scalar} on
the same interval (Dormand--Prince 5(4), or dense Gauss--Kronrod
on~\eqref{eq:p1stable}); expect agreement to $\approx 10^{-12}$ relative, in the
style of \texttt{test-transit-chain.R}'s RK cross-checks.

% ===========================================================================
\section{P2 --- two or more slow ligands}
\label{sec:p2}

\begin{correction}{The multi-slow-ligand advance is linear, indeed compartmental --- not nonlinear}
With buprenorphine plus a long-residence antagonist, both keep differential
states and there is \textbf{no scalar reduction} --- but the coupled system is
\textbf{linear time-varying}, not nonlinear. The product
$\kon[,i]\Cesub{i}(t)\rhoz$ is linear in the occupancies because $\Cesub{i}(t)$ is a
known input, not a state; the only ``nonlinear-looking'' factor,
$\rhoz=1-\sum\rho$, is itself linear in the states.

Carry the free fraction explicitly and the structure is a genuine
\textbf{compartmental (Markov) generator}. For slow ligands with coefficients
$a_i(t)=\kon[,i]\Cesub{i}(t)/(1+\Phi(t))$, the states $(\rho_1,\dots,\rho_N,\rhoz)$
obey $\dot\rho=K(t)\rho$ with, for two slow ligands $B$ and $A$,
%
\begin{equation}
\dd{}{t}
\begin{pmatrix}\rho_B\\[2pt]\rho_A\\[2pt]\rhoz\end{pmatrix}
=
\begin{pmatrix}
-\koff[,B] & 0 & a_B\\[2pt]
0 & -\koff[,A] & a_A\\[2pt]
\koff[,B] & \koff[,A] & -(a_B+a_A)
\end{pmatrix}
\begin{pmatrix}\rho_B\\[2pt]\rho_A\\[2pt]\rhoz\end{pmatrix}.
\label{eq:generator}
\end{equation}
%
$K(t)$ is \textbf{Metzler} (off-diagonals $\ge0$) with \textbf{zero column
sums} (conservative) --- a continuous-time Markov generator. Hence
$\exp(K\Delta)$ is a nonnegative, column-stochastic matrix: a matrix-exponential
step keeps every occupancy $\ge0$ and $\sum\rho=1$ \textbf{exactly}, so the
physical simplex is preserved by construction, with no clipping. (The negative
off-diagonals one would see in a 2-state $(\rho_B,\rho_A)$ reduction are an
artifact of eliminating $\rhoz$; keep $\rhoz$ and they vanish.) The fast ligands
then redistribute $\rhoz$'s share by~\eqref{eq:qss} at each output time.
\end{correction}

\paragraph{The step.} Over $[0,h]$, $K(t)$ varies through $a_i(t)$ (each a
M\"obius function of $t$), so the exact propagator is a first-order
\textbf{Magnus} exponential $\exp\!\bigl(\int_0^h K\,\mathrm{d}t\bigr)$. Its
entries are the same elementary $\int a_i\,\mathrm{d}t$ logarithms as
\S\ref{sec:p1closed}, so the integral is available in closed form; use a
midpoint / averaged $\bar K$ for second-order accuracy, or sub-step. Advance
with the \textbf{augmented unit-state} device already in
\texttt{advanceTransitChain()}. The only approximation is the within-step
time-variation of $K$ --- identical in spirit to \texttt{advanceClosedFormGas()}
holding its coupling at the start-of-step value --- and it is exact to rounding
when $K$ is constant on the step. So ``no closed form survives'' is true only of
a \emph{scalar} closed form; there is an exact \textbf{matrix}-exponential
advance per interval, well conditioned and simplex-preserving. For $N$ slow
ligands it is an $(N{+}1)$-state generator.

% ===========================================================================
\section{P4 --- backward compatibility}
\label{sec:p4}

Calibrate each fast agonist so $C_{50,j}=\mathrm{MEAC}_j$. Then
$\phi_j=\Cesub{j}/C_{50,j}=\Cesub{j}/\mathrm{MEAC}_j$, so the crowding factor is
\textbf{exactly the legacy opioid scale}, $\Phi=\sum_{\mathrm{fast}}
\Cesub{j}/\mathrm{MEAC}_j=U$. With \textbf{no slow ligand} present,
%
\begin{equation}
\sum_{\mathrm{fast}}\rho_j=\frac{U}{1+U},\qquad \rhoz=\frac{1}{1+U}.
\label{eq:p4base}
\end{equation}

\subsection{The identity, and what ``unchanged'' means}
\eqref{eq:p4base} is exact, but $U/(1+U)\ne U$ --- it is a \textbf{saturating
transform} of $U$. So ``leaves the additive panels unchanged'' holds in the
\textbf{monotone / order-preserving} sense, not as numerical equality. One
consequence is desirable: $U=1$ (one MEAC of a single agonist) gives
$\sum\rho=\tfrac12$, i.e.\ \textbf{MEAC $=$ the 50\%-occupancy point} ---
arguably more defensible than a bare linear sum, and consistent with reading
MEAC as a half-effect concentration.

\begin{correction}{P4 has an exact invariant, not merely a bounded deviation}
If the legacy panel is read not as raw occupancy but as the fast-ligand
\textbf{odds} (fast-bound over free), the correspondence is exact \emph{and}
survives the presence of a slow ligand. With $\sum_{\mathrm{slow}}\rho$
arbitrary,
$\rhoz=(1-\sum_{\mathrm{slow}}\rho)/(1+\Phi)$ and
$\sum_{\mathrm{fast}}\rho=(1-\sum_{\mathrm{slow}}\rho)\,\Phi/(1+\Phi)$, so
%
\begin{equation}
\frac{\sum_{\mathrm{fast}}\rho}{\rhoz}=\Phi=U
\qquad\text{exactly, independent of }\textstyle\sum_{\mathrm{slow}}\rho.
\label{eq:invariant}
\end{equation}
%
The fast-ligand occupancy odds equals $U$ whether or not buprenorphine is
present. This is the real backward-compatibility theorem: \textbf{report
$U=\sum_{\mathrm{fast}}\rho/\rhoz$ and the legacy scale is reproduced
identically}, with buprenorphine's effect entering only through the separate
occupancy term it adds, never by corrupting $U$.
\end{correction}

\subsection{The deviation bound (for a panel that insists on raw occupancy)}
If instead the panel plots raw fast occupancy $\sum_{\mathrm{fast}}\rho$, the
deviation from the no-slow value is bounded by buprenorphine's own occupancy:
%
\begin{equation}
\Bigl|\,\textstyle\sum_{\mathrm{fast}}\rho-\frac{U}{1+U}\,\Bigr|
=\sum_{\mathrm{slow}}\rho\cdot\frac{U}{1+U}\le\sum_{\mathrm{slow}}\rho,
\end{equation}
%
$\to 0$ as $\sum_{\mathrm{slow}}\rho\to0$. Either way the additive panels are
untouched until a partial agonist or antagonist appears.

% ===========================================================================
\section{P3 --- operational transduction and the analgesia/ventilation asymmetry}

Replace the linear form \eqref{eq:effect} with an operational model (Black \&
Leff 1983). Give each ligand an intrinsic efficacy $\tau_i$ and each
\textbf{endpoint} a reserve / gain $\gamma$; with transduction Hill $n$,
%
\begin{equation}
\frac{E}{E_{\max}}=\frac{\Theta^{\,n}}{1+\Theta^{\,n}},\qquad
\Theta\equiv\gamma\sum_i\tau_i\rho_i.
\label{eq:operational}
\end{equation}

\begin{correction}{The low-reserve limit recovers Olofsen's \eqref{eq:effect} only if $n=1$}
For a low-reserve endpoint $\Theta\ll1$, so
$\Theta^{n}/(1+\Theta^{n})\approx\Theta^{n}$. This is \textbf{linear in
occupancy only when $n=1$}: then $E/E_{\max}\approx\gamma\sum\tau_i\rho_i$, and
writing ventilation as depression $E=E_0(1-\sum\alpha_i\rho_i)$ identifies
$\alpha_i=\gamma_{\mathrm{vent}}\tau_i$. For $n\ne1$ the low-reserve limit is a
power law, not \eqref{eq:effect}. State the $n=1$ condition whenever the
reduction is invoked.
\end{correction}

\paragraph{One receptor, different reserve per endpoint, reproduces the
asymmetry.} A single $\mu$ model with a larger reserve for \textbf{analgesia}
than for \textbf{ventilation} gives buprenorphine's clinical split:
\begin{itemize}[nosep,leftmargin=1.4em]
\item \textbf{Analgesia} (high $\gamma$): $\Theta$ saturates
$\Theta^n/(1+\Theta^n)\to1$ even at the modest $\sum\tau_i\rho_i$ a low-efficacy
agonist produces $\Rightarrow$ buprenorphine looks \textbf{near-full}.
\item \textbf{Ventilation} (low $\gamma$): $E/E_{\max}\approx\gamma\sum\tau_i
\rho_i$ stays sub-saturating, so a low-$\tau$ ligand gives a \textbf{ceiling}
$\approx\gamma_{\mathrm{vent}}\tau_{\mathrm{bup}}<1$ $\Rightarrow$ the partial
agonist.
\end{itemize}
This is the receptor-reserve account (standard for the $\mu$
ventilation/analgesia asymmetry) and needs neither NOP nor biased signalling ---
appropriate to the stated $\mu$-only scope.

\paragraph{Mapping Olofsen's $\alpha$, and the identifiability caveat.} Preserve
the ventilation fit by setting $\gamma_{\mathrm{vent}}\tau_i=\alpha_i$
(Olofsen's values) in the $n=1$ low-reserve limit. The efficacy \emph{ratio}
$\tau_{\mathrm{bup}}/\tau_F=\alpha_{\mathrm{bup}}/\alpha_F$ is then
endpoint-independent; the absolute ceiling differs by endpoint only through
$\gamma$. \textbf{Caveat:} in a single endpoint $\tau$ and $\gamma$ enter only
as the product $\gamma\tau$, so ventilation data alone fix
$\gamma_{\mathrm{vent}}\tau_i$ but not the split. The analgesia reserve
$\gamma_{\mathrm{analgesia}}$ (equivalently
$\gamma_{\mathrm{analgesia}}/\gamma_{\mathrm{vent}}$) is a \textbf{second
parameter} that must come from analgesia data or be set to reproduce
``near-full agonist for analgesia''; it cannot be read off the ventilation fit.

% ===========================================================================
\section{P5 --- tolerance as a one-parameter manifold}

Olofsen gives two states (naive, chronic) differing in
$(\kon[,B],\alpha_B,C_{50,F})$ with $\koff[,B]$ fixed. The app exposes four
tolerance levels. With only two anchors, \textbf{only a one-parameter monotone
path is identifiable} --- no data support curvature, so interpolate linearly in
the natural transform of each parameter across $\theta\in[0,1]$, anchored at
$\theta=0$ (naive) and $\theta=1$ (chronic):
\begin{itemize}[nosep,leftmargin=1.4em]
\item $\log\kon[,B]$ \textbf{linear in $\theta$} ($0.26\to0.085$, a $\sim3\times$
range).
\item $\log C_{50,F}$ \textbf{linear in $\theta$} ($0.68\to2.23$, a
$\sim3.3\times$ range).
\item $\alpha_B$ on a \textbf{bounded scale} --- linear, or logit
$\log(\alpha/(1-\alpha))$ linear in $\theta$ to keep $\alpha\in(0,1)$
($0.81\to0.48$).
\item $\koff[,B]$ \textbf{fixed} ($0.019$), as given.
\end{itemize}
Four levels then sit at $\theta=0,\tfrac13,\tfrac23,1$. Anything richer than
monotone-in-transform over-fits two points.

% ===========================================================================
\section{Validation statements}

\begin{table}[h]
\centering
\begin{tabular}{@{}l p{0.42\linewidth} p{0.26\linewidth}@{}}
\toprule
Result & Checked against & Expected \\
\midrule
P1 advance \eqref{eq:p1stable} & adaptive DP5(4) / Gauss--Kronrod on \eqref{eq:scalar} & $\approx10^{-12}$ rel. \\
P1 degenerate limits \S\ref{sec:p1deg} & decay / constant-coeff. / erf forms & exact \\
P2 matrix step \S\ref{sec:p2} & dense RK of $(N{+}1)$-state system; $\sum\rho=1$, $\rho\ge0$ & $\approx10^{-10}$, simplex exact \\
P3 reduction (Corr.\ in \S9) & linear form \eqref{eq:effect} at $n=1$, small $\Theta$ & matches slope \\
P4 invariant \eqref{eq:invariant} & algebra; simulation with/without slow ligand & exact \\
\bottomrule
\end{tabular}
\end{table}

\paragraph{The Nasser 2014 occupancy check is loose, not tight.} Nasser et al.\
2014 give buprenorphine whole-brain $\mu$ occupancy $=91.4\,C/(0.67+C)$ percent
vs plasma concentration, offered as a check on $K_{D,B}$. The ventilation-fit
value is $K_{D,B}=\koff/\kon=0.019/0.26\approx0.073~\mathrm{ng\,mL^{-1}}$
(naive), $0.019/0.085\approx0.22$ (chronic); the occupancy half-max is at
$\mathbf{0.67~ng\,mL^{-1}}$, about $3$--$9\times$ higher. The gap is expectable
--- plasma vs biophase, PET whole-brain vs ventilatory biophase, and $\kon$
fitted in effect-site units --- but it makes Nasser an \textbf{order-of-magnitude
sanity check, not a tight validation} of $K_{D,B}$. Do not present the two as
agreeing closely.

% ===========================================================================
\section{Where an implementation would sit}

The engine is a \textbf{pharmacodynamic post-process}, so it would run after
\texttt{simCpCe()} has produced every opioid's $\Cesub{i}(t)$, not inside the PK
engines --- analogous to \texttt{applyAntiseizure\-\allowbreak Interactions()}
running after each drug is simulated alone. It would:
\begin{enumerate}[nosep,leftmargin=1.6em]
\item partition the $\mu$ ligands into slow (differential state) and fast
(conditional equilibrium \eqref{eq:qss}) by $\koff$;
\item advance the slow occupancies interval by interval on the shared,
non-uniform time line --- \S6 (one slow ligand) or \S\ref{sec:p2} (several);
\item form $\Phi(t)$ and the fast occupancies by \eqref{eq:qss} at each output
time;
\item map occupancy to effect by \eqref{eq:effect} now, or the operational form
\eqref{eq:operational} if and when an analgesia endpoint is added.
\end{enumerate}
It must leave the additive opioid panels identical until buprenorphine or an
antagonist is dosed --- guaranteed exactly by \eqref{eq:invariant}, and to
within $\sum_{\mathrm{slow}}\rho$ by raw occupancy.

\paragraph{Open for Shafer --- decisions this note does not make.}
\begin{itemize}[nosep,leftmargin=1.4em]
\item \textbf{[decision]} whether the production P1 step uses the
incomplete-gamma closed form (\S\ref{sec:p1closed}) or the bounded-kernel
quadrature (\S\ref{sec:p1overflow}). This note recommends the quadrature with
the closed form as validator; it is a decision because it sets what the tests
pin.
\item \textbf{[decision]} the analgesia reserve $\gamma_{\mathrm{analgesia}}$ ---
a second endpoint's worth of information that ventilation data cannot supply.
Needed before P3 can draw an analgesia curve; until then the engine is
ventilation-only and \eqref{eq:effect} stands.
\item \textbf{[decision]} whether the four tolerance levels are
$\theta=0,\tfrac13,\tfrac23,1$ or another spacing.
\end{itemize}

\end{document}
